\section{Conclusions}
\label{sec:conclusion}

TRACT counts white-tailed deer from thermal road transects, and decomposing its error moves
the problem rather than solving it. What a practitioner would try to buy with a better
detector is lost after detection, in the stage that decides whether a track is a new
animal. Every improvement to detection and association that we tried left the count no
better or worse, and learned confirmers of three capacities lost to a three-parameter rule.
The reason is supervision rather than model class. The confirmation stage sees roughly two
hundred positive tracks, and duplicates of animals already counted sit \emph{between} the
positives and the false candidates on every feature we measured. What separates a first
sighting from a re-sighting is therefore precisely what a confirmer cannot learn at this
scale. Two boundaries follow for anyone adopting the result. The pipeline does not transfer
unaltered: a detector trained without one site and its survey night scores it low enough
that an absolute confidence threshold can reject every track (\cref{sec:loso}). And the
under-count is an operating-point choice rather than a limit, since at the threshold where
signed bias crosses zero the same pipeline recovers 81 of the 83 held-out animals in
aggregate (\cref{sec:operating}). We release the dataset, per-individual annotations,
evaluation harness and decomposition tooling. The gap between what is reachable and what is
counted is large, precisely localized and entirely open.
